%=============================================
% DIAPOSITIVA FINAL UNIFICADA
% Kevin
%=============================================

\begin{frame}{Conclusiones, Optimización y Validación del Modelo}

    \begin{columns}[T]

        %=========================================
        % COLUMNA IZQUIERDA
        %=========================================
        \column{0.48\textwidth}

        \begin{block}{Significado del Mínimo}
            \scriptsize
            La función de error depende de cuatro variables:
            \[
                E(x,y,z,A_0)
            \]

            El objetivo fue encontrar el punto donde el error fuera mínimo:
            \[
                m^* = \arg\min_{m \in \mathbb{R}^4} E(m)
            \]

            Ese mínimo representa la ubicación espacial que mejor reproduce las amplitudes observadas por los sensores.
        \end{block}

        \vspace{0.15cm}

        \begin{exampleblock}{Optimización con Gauss-Newton}
            \scriptsize
            \begin{itemize}
                \item Identificó automáticamente la región de menor error.
                \item Evitó evaluar infinitas combinaciones manualmente.
                \item Convergió en solo \textbf{8 iteraciones}.
                \item Precisión obtenida superior al \textbf{99.7\%}.
            \end{itemize}
        \end{exampleblock}

        \vspace{0.15cm}

        \begin{block}{Conclusión Principal}
            \scriptsize
            Resolver el mínimo de la función de error equivale físicamente a localizar la fuente sísmica.
        \end{block}

        %=========================================
        % COLUMNA DERECHA
        %=========================================
        \column{0.48\textwidth}

        \begin{block}{Validación del Modelo}
            \centering
            \includegraphics[width=0.95\textwidth]{Imagenes/superficie_error_3d.png}
        \end{block}

        \vspace{0.1cm}

        \begin{block}{Limitaciones y Mejoras Futuras}
            \scriptsize
            \begin{itemize}
                \item Se asumió un medio homogéneo.
                \item El ruido afecta la precisión de las mediciones.
                \item Más sensores mejorarían la resolución espacial.
                \item Futuramente podrían incorporarse modelos físicos más complejos.
            \end{itemize}
        \end{block}

    \end{columns}

    \vspace{0.2cm}

    \begin{center}
        \footnotesize
        \textbf{
        Este proyecto integró cálculo multivariado, geometría espacial y optimización
        para resolver un problema inverso de localización sísmica.
        }
    \end{center}

    \vspace{0.1cm}

    \begin{center}
        \small
        \textcolor{blue}{\textbf{¡Gracias por su atención!}}
    \end{center}

\end{frame}